% !PRINTED 374
\weberpage{374}
%<<CONT
eine lineare Substitution, und
\[
\begin{pmatrix}m, & \mu\\ n, & \nu\end{pmatrix} = \begin{pmatrix}\alpha, & \beta\\ \gamma, & \delta\end{pmatrix}\begin{pmatrix}m', & \mu'\\ n', & \nu'\end{pmatrix},
\]
woraus folgt:
\[
m = \alpha m' + \beta n', \qquad n = \gamma m' + \delta n',
\]
also
\[
\frac{m}{n} = \dfrac{\alpha\dfrac{m'}{n'}+\beta}{\gamma\dfrac{m'}{n'}+\delta},
\]
was zu beweisen war. Die Aequivalenz ist eine eigentliche oder eine uneigentliche, je nachdem $\varepsilon = \varepsilon'$ oder $\varepsilon = -\varepsilon'$ ist.

Jede Zahl ist mit ihrer entgegengesetzten und mit ihrer reciproken uneigentlich äquivalent, denn es ist
\[
-x = \begin{pmatrix}-1, & 0\\ 0, & 1\end{pmatrix}(x), \qquad \frac{1}{x} = \begin{pmatrix}0, & 1\\ 1, & 0\end{pmatrix}(x).
\]

\begin{center}
\S.\,121.

\gesperrt{Entwickelung äquivalenter Zahlen in Kettenbrüchen.}
\end{center}

Wenn von zwei äquivalenten Zahlen die eine irrational ist, so ist es die andere auch, und wenn die eine reell ist, so ist es auch die andere. Wir machen diese beiden Annahmen, setzen also in
\weq{1}{y = \dfrac{\alpha x + \beta}{\gamma x + \delta}, \qquad \alpha\delta - \beta\gamma = \varepsilon}
$x$ als irrational und reell voraus, und wollen nun untersuchen, wie sich aus der Kettenbruchentwickelung von $x$ die Kettenbruchentwickelung von $y$ herleiten lässt.

Es sei also $x$ in einen Kettenbruch mit der Schlusszahl $x_n$ entwickelt:
\weq{2}{x = (a_0, a_1, a_2 \ldots a_{n-1}, x_n),}
worin wir $n$ vorläufig unbestimmt lassen.

Durch die Näherungsbrüche ausgedrückt, wird
\weq{3}{x = \dfrac{P_n x_n + P_{n-1}}{Q_n x_n + Q_{n-1}}, \qquad P_n Q_{n-1} - Q_n P_{n-1} = (-1)^n.}
